\item Suponga que un átomo de hidrógeno se encuentra en el estado cuántico $2s$, con una función de onda radial dada por:
$$ \psi_{2s}(r) = \frac{1}{4\sqrt{2\pi a_0^3}} \left( 2 - \frac{r}{a_0} \right) e^{-r/2a_0} $$
Si el electrón se ubica a $r = a_0$, calcule:
\begin{enumerate}
    \item el valor de la función de onda $\psi_{2s}(a_0)$,
    \item la densidad de probabilidad $\left|\psi_{2s}(a_0)\right|^2$, y
    \item la densidad de probabilidad radial $P_{2s}(a_0)$.
\end{enumerate}